\begin{textsample}{POS Dim 3 – ms – Score 7.00 – ms/ms\_inf\_e\_ef\_27.txt}  \label{ex:f3_pos_239}
7.3.1 Objetivos de Aprendizagem e Desenvolvimento Tendo em vista as especificidades da Educação \textbf{Infantil}, as \textbf{crianças} precisam vivenciar experiências significativas na perspectiva de que essas possam incidir sobre a aprendizagem e o desenvolvimento delas.
Na BNCC, reconhecendo as características das diferentes faixas etárias, os objetivos de aprendizagem e desenvolvimento foram organizados em três grupos : \textbf{bebês} ( de 0 – 1 ano e 6 \textbf{meses} ), \textbf{crianças} bem \textbf{pequenas} ( de 1 ano e 7 \textbf{meses} – 3 anos e 11 \textbf{meses} ) e \textbf{crianças} \textbf{pequenas} ( de 4 anos – 5 anos e 11 \textbf{meses} ), porém não inviabilizam outras formas de organização nos municípios, desde que as propostas respeitem o tempo e o modo de aprender das \textbf{crianças}.
Assim, no presente documento, além dos objetivos definidos pela BNCC, foram inseridos outros como forma de ampliar os conhecimentos das \textbf{crianças}, cujas preservação, divulgação e ampliação são papéis das instituições educativas.
O desafio de construir coletivamente um currículo para a Educação \textbf{Infantil} que respeite a pluralidade e diversidade das \textbf{crianças} sul-mato-grossenses, coloca em pauta outros fatores que também devem ser considerados na consolidação de um trabalho pedagógico de qualidade : a infraestrutura, garantindo espaços diferenciados, materiais e \textbf{brinquedos} em quantidade suficiente ; a proporção adulto/criança ; a gestão do tempo nas instituições ; a formação inicial e continuada dos profissionais : professores, gestores, coordenadores, assim como os administrativos, para que esses entendam a especificidade da faixa etária ; e, também, a interação da instituição com as famílias e outras instituições responsáveis pela educação das \textbf{crianças} com vistas a subsidiar os projetos educativos em parceria com a comunidade. 7.4 Os Campos de Experiências e as Ações Didáticas O cuidar e o educar são princípios fundamentais na Educação Básica.
Para as ações didáticas, o professor deve levar em consideração a relevância social e cultural dos objetivos, assegurando a progressão nos saberes para a formação integral das \textbf{crianças}, nas dimensões afetivas, cognitivas, físicas, sociais e culturais.
Nesse sentido, as ações didáticas aqui propostas são possibilidades que poderão ser ampliadas nas instituições a partir de concepções comuns de educação, \textbf{criança}, ensino, aprendizagem e avaliação, definidas em cada Proposta Pedagógica.
Nessas escolhas coletivas o professor terá como imprimir a intencionalidade pedagógica, planejando de forma sistemática, porém de maneira contextualizada visando à garantia do direito das \textbf{crianças} de ampliar seus conhecimentos e de frequentar uma instituição educativa de qualidade.

% matched lemmas: bebê, brinquedo, criança, infantil, mês, pequeno
\end{textsample}
